\BourbakiExercises{习题}

\begin{enumerate}
\def\labelenumi{\arabic{enumi})}
\item
  设 A 是一个半单环，M 是一个有限生成的 A-模。对 M 的任一子模 N，我们把 N 在 \(M^{*}\) 中的正交记作 \(N'\)。证明 \(N'\) 在 M 中的正交等于 N，且我们有 \(\text{long}_{\text{A}}(\text{N}) + \text{long}_{\text{A}}(\text{N}') = \text{long}_{\text{A}}(\text{M})\)。映射 \(N \mapsto N'\) 是从 M 的子模所成的集到 \(M^{*}\) 的子模所成的集上的一个双射。若 L 是 M 的另一个子模，则 \(L \cap N\) 的正交是 \(L' + N'\)。
\item
  设 A 是一个环。证明 \((A,A)\)-双模 \(_{s}A_{d}\) 是半单的，当且仅当 A 同构于有限多个拟单环的一个乘积。
\item
  我们称环的一个（左或右）理想 \(\mathfrak{a}\) 是幂零的，如果存在一个整数 \(n\)，使得由 \(n\) 个元素（属于 \(\mathfrak{a}\)）构成的每个乘积都为零。
\end{enumerate}

\begin{enumerate}
\def\labelenumi{\alph{enumi})}
\item
  设 \(\mathfrak{a}\) 是 A 的一个幂零左理想。证明双边理想 \(\mathfrak{aA}\)（由 \(\mathfrak{a}\) 生成）是幂零的。
\item
  我们称环 A 是半素的，如果它不含任何非零的幂零双边理想。我们称 A 的一个双边理想 \(\mathfrak{a}\) 是半素的，如果商环 A/ \(\mathfrak{a}\) 是半素的。证明一族半素（双边）理想的交是一个半素理想，且一个半素理想包含 A 的每个幂零理想。
\item
  一个素（双边）理想（VIII, 第 133 页，习题 15）是半素的。一个双边理想 a 是半素的，当且仅当它是包含它的诸素理想的交。（若 \(a \notin a\)，考虑一个在包含 a 的双边理想之中极大、且不包含 a 的任何幂的理想。）
\end{enumerate}

¶ 4) 设 A 是一个半素左诺特环（习题 3），S 是其可消去元素所成的集（参见 VIII, 第 75 页，习题 8）。

\begin{enumerate}
\def\labelenumi{\alph{enumi})}
\item
  证明 A 的每个非零（左或右）理想都包含一个非幂零的元素。（设 a 是 A 的一个由幂零元素组成的右理想。取 a 的一个非零元素 x，使其左零化子极大。注意，对 A 的每个元素 a 及每个整数 \(n \geqslant 1\)（满足 \((xa)^{n} \neq 0\)），\((xa)^{n}\) 的左零化子等于 x 的左零化子，并由此推出我们有 xax = 0。）
\item
  证明 S 在 A 上允许一个左分式演算（VIII, 第 18 页，习题 18）。（对 \(a \in \mathrm{A}\) 与 \(s \in \mathrm{S}\)，证明由 A 中元素 \(x\)（满足 \(xa \in \mathrm{As}\)）组成的集是 A 的一个不可约左理想（VIII, 第 75 页，习题 7），并利用 \(a\) ) 与 VIII, 第 75 页的习题 8, \(d\) ) 得出结论。）
\end{enumerate}

\(c\) ) 证明环 \(\mathrm{S}^{-1}\mathrm{A}\) 是半单的（``戈尔迪定理'' (1)）。（设 \(\mathfrak{b}\) 是 \(\mathrm{S}^{-1}\mathrm{A}\) 的一个不可约左理想。证明 \(\mathrm{A} \cap \mathfrak{b}\) 是 A 的一个不可约左理想，然后利用 VIII, 第 75 页的习题 8, d) 证明 \(\mathfrak{b} = \mathrm{S}^{-1}\mathrm{A}\)。再利用 VIII, 第 75 页的习题 7, b) 得出结论。）

\( ^{(1)} \) 关于此结果的其他处理方式，参见 X, §1, 第 171 页，习题 25 及《交换代数》II, §2, 第 135 页，习题 26。

\begin{enumerate}
\def\labelenumi{\arabic{enumi})}
\setcounter{enumi}{4}
\tightlist
\item
  \begin{enumerate}
  \def\labelenumii{\alph{enumii})}
  \tightlist
  \item
    设 A 是一个半素环，使得每族单生成左理想都有一个极小元素。证明 A 是半单的。（利用 VIII, 第 130 页的习题 7，构造 A 中幂等元的序列 \((e_n)\) 与 \((f_n)\)，使得诸理想 \(Ae_n\) 极小，且我们有 \(A f_{n-1} = Ae_n \oplus Af_n\)，并证明 \(f_n\) 当 \(n\) 充分大时为零。）
  \end{enumerate}
\end{enumerate}

\begin{enumerate}
\def\labelenumi{\alph{enumi})}
\setcounter{enumi}{1}
\tightlist
\item
  设 K 是一个交换域，V 是 K 上的一个无穷维右向量空间，A 是 \(\operatorname{End}_{\mathrm{K}}(\mathrm{V})\) 中由 \(1_{V}\) 与有限秩自同态生成的 K-子代数。证明 A 的每个左理想都包含一个极小理想。
\end{enumerate}

\begin{enumerate}
\def\labelenumi{\arabic{enumi})}
\setcounter{enumi}{5}
\tightlist
\item
  设 A 是交换域 K 上有限秩的半单交换代数，设 \(K_{i}\)（对 \(i \in I\)）是它的单因子，它们是 K 的有限扩张。设 B 是 A 的一个子代数，\(E_{j}\)（\(j \in J\)）是它的单因子。证明存在 I 的两两不相交的非空子集 \((\Lambda_{j})_{j \in J}\)，使得 \(E_{j}\) 单射到 \(\prod_{i \in \Lambda_{j}} K_{i}\) 中（对 \(j \in J\)）。对每个 \(i \in \Lambda_{j}\)，\(pr_{i}\) 诱导一个从 \(E_{j}\) 到子域 \(E_{ij}\)（\(K_{i}\) 的）的同构。研究其逆命题。推出：若每个 \(K_{i}\) 在 K 上可分，则 A 只含有限多个子代数（参见 V, §6, No.~4, 第 30 页，命题 3）。
\end{enumerate}

